% from ega4-18-fr.tex
\oldpage[IV-4]{190}

\begin{definition}[19.2.1]
\EGAReaderTarget{ega4.c04.def.19-2-1}{IV.19.2.1}
Soient $u:X\to S$ un morphisme de préschémas, $\mathscr F$ un $\sh{O}_X$-Module quasi-cohérent. On dit qu'une suite $(f_i)_{1\leq i\leq n}$ de sections de $\mathscr F$ au-dessus de $X$ est \emph{transversalement $\mathscr F$-régulière relativement à $S$} (ou relativement à $u$) si la suite $(f_i)$ est $\mathscr F$-régulière et telle que les $\sh{O}_X$-Modules
\[
  \mathscr F\big/\Bigl(\sum_{j\leq i}f_j\mathscr F\Bigr)
\]
soient $S$-plats pour $0\leq i\leq n$. On dit qu'un Idéal quasi-cohérent $\mathscr J$ de $\sh{O}_X$ est \emph{transversalement régulier relativement à $S$} (ou relativement à $u$) en un point $x\in\operatorname{Supp}(\sh{O}_X/\mathscr J)$ s'il existe un voisinage ouvert $U$ de $x$ et une suite finie $(f_i)$ de sections de $\mathscr J$ au-dessus de $U$ qui est transversalement $\sh{O}_U$-régulière relativement à $S$ et qui engendre $\mathscr J|U$.

Dans une $A$-algèbre $B$, on dit qu'un idéal $\mathfrak J$ est \emph{transversalement régulier relativement à $A$} si $\widetilde{\mathfrak J}$ est transversalement régulier relativement à $S=\Spec(A)$ en tout point de $V(\mathfrak J)$ dans $\Spec(B)$.
\end{definition}

\begin{definition}[19.2.2]
\EGAReaderTarget{ega4.c04.def.19-2-2}{IV.19.2.2}
Soient $S$ un préschéma, $X$, $Y$ deux $S$-préschémas, $u:Y\to X$ un $S$-morphisme qui est une immersion. On dit que $u$ est une \emph{immersion transversalement régulière relativement à $S$ en un point $y\in Y$} s'il existe un voisinage $V$ de $u(y)$ dans $X$ tel que le sous-préschéma induit sur $u(Y)\cap V$ par le sous-préschéma de $X$ associé à $u$ soit défini par un Idéal de $\sh{O}_U$ transversalement régulier relativement à $S$ au point $u(y)$.

Il est clair que l'ensemble des points de $Y$ où une immersion est transversalement régulière relativement à $S$ est \emph{ouvert} dans $Y$; s'il est égal à $Y$, on dit simplement que $u$ est une \emph{immersion transversalement régulière relativement à $S$}.

Si $S\to T$ est un morphisme \emph{plat}, alors, si $\mathscr J$ est un Idéal de $\sh{O}_X$ transversalement régulier relativement à $S$ en $x\in X$, il l'est aussi relativement à $T$ ($0_{\mathrm I}$, 6.2.1). Une $S$-immersion $Y\to X$ transversalement régulière en un point $y\in Y$ relativement à $S$ est donc aussi une $T$-immersion transversalement régulière relativement à $T$ en ce point.
\end{definition}

\begin{proposition}[19.2.3]
\label{IV.19.2.3}
Soient $S$ un préschéma, $X$, $Y$ deux $S$-préschémas, $u:Y\to X$ un $S$-morphisme qui est une immersion transversalement régulière relativement à $S$ en un point $y\in Y$. Alors, pour tout changement de base $S'\to S$, si l'on pose $X'=X\times_S S'$, $Y'=Y\times_S S'$, l'immersion $u'=u_{(S')}:Y'\to X'$ est transversalement régulière relativement à $S'$ en tout point $y'\in Y'$ au-dessus de $y$.
\end{proposition}

Cela résulte aussitôt de la définition et de ($0$, 15.1.15).

\begin{proposition}[19.2.4]
\label{IV.19.2.4}
Soient $S$ un préschéma, $X$, $Y$ deux $S$-préschémas, $u:Y\to X$ un $S$-morphisme qui est une immersion. On suppose, soit que $S$ et $X$ sont localement noethériens, soit que les morphismes structuraux $g:X\to S$, $h:Y\to S$ soient localement de présentation finie. Soient $y$ un point de $Y$, $s$ son image dans $S$. Les conditions suivantes sont équivalentes:
\begin{enumerate}
  \item[a)] $u$ est transversalement régulière relativement à $S$ au point $y$.
  \item[b)] Les morphismes $g$ et $h$ sont plats au voisinage des points $u(y)$ et $y$ respectivement, et si l'on pose $X_s=g^{-1}(s)$, $Y_s=h^{-1}(s)$, l'immersion $u_s:Y_s\to X_s$ est régulière au point $y$.
  \item[b$'$)] Les morphismes $g$ et $h$ sont plats au voisinage des points $u(y)$ et $y$ respectivement, et, pour tout changement de base $S'\to S$, si l'on pose $X'=X\times_S S'$, $Y'=Y\times_S S'$, l'immersion $u'=u_{(S')}:Y'\to X'$ est régulière en tout point de $Y'$ au-dessus de $y$.
  \item[c)] Le morphisme $h$ est plat au voisinage du point $y$, et l'immersion $u$ est régulière au voisinage du point $y$.

